\documentclass[12pt]{article}
\usepackage{amsmath,amssymb,amsfonts}
\usepackage[margin=1in]{geometry}
\usepackage{bm}

\newcommand{\vect}[1]{\mathbf{#1}}

\begin{document}

\section*{Problem 12}

\textbf{Problem:} A second-rank antisymmetric tensor $F_{\mu\nu}$ is characterized by three independent quantities. Consequently, a vector $H_\mu$ can be associated with such a tensor. Prove that $H_\mu = \frac{1}{2}\varepsilon_{\mu\alpha\beta}F_{\alpha\beta}$ if and only if $F_{\alpha\beta} = \varepsilon_{\alpha\beta\mu}H_\mu$.

\textit{(Note: This problem and Problem 11 are stated equivalently; the solution from Problem 11 applies. Here we verify the relationship explicitly with components.)}

\bigskip
\textbf{Solution:}

The general proof was given in Problem 11. Here we verify the correspondence component by component.

\subsection*{Explicit verification}

An antisymmetric $3\times 3$ tensor $F_{\mu\nu}$ ($F_{\mu\nu} = -F_{\nu\mu}$) has three independent components. Following the book's convention (Eq.~1.96), we write:
\[
F_{\mu\nu} = \begin{pmatrix} 0 & F_{12} & F_{13} \\ -F_{12} & 0 & F_{23} \\ -F_{13} & -F_{23} & 0 \end{pmatrix}.
\]

The associated vector $H_\mu = \frac{1}{2}\varepsilon_{\mu\nu\lambda}F_{\nu\lambda}$ gives:
\begin{align}
H_1 &= \tfrac{1}{2}\varepsilon_{1\nu\lambda}F_{\nu\lambda} = \tfrac{1}{2}(\varepsilon_{123}F_{23} + \varepsilon_{132}F_{32}) = \tfrac{1}{2}(F_{23} - F_{32}) = F_{23}, \\
H_2 &= \tfrac{1}{2}\varepsilon_{2\nu\lambda}F_{\nu\lambda} = \tfrac{1}{2}(\varepsilon_{213}F_{13} + \varepsilon_{231}F_{31}) = \tfrac{1}{2}(-F_{13} + F_{13}) \nonumber \\
    &= \tfrac{1}{2}(\varepsilon_{231}F_{31} + \varepsilon_{213}F_{13}) = \tfrac{1}{2}(F_{31} + (-1)F_{13}) = \tfrac{1}{2}(-F_{13} - F_{13}) \nonumber \\
    &\quad\text{Wait, let us be more careful:} \nonumber \\
H_2 &= \tfrac{1}{2}\sum_{\nu,\lambda}\varepsilon_{2\nu\lambda}F_{\nu\lambda}. \nonumber
\end{align}

The nonzero terms of $\varepsilon_{2\nu\lambda}$ are: $\varepsilon_{213} = -1$ and $\varepsilon_{231} = +1$. So:
\begin{align}
H_2 &= \tfrac{1}{2}[\varepsilon_{213}F_{13} + \varepsilon_{231}F_{31}] = \tfrac{1}{2}[(-1)F_{13} + (1)(-F_{13})] = \tfrac{1}{2}(-2F_{13}) = -F_{13} = F_{31}.
\end{align}

Similarly:
\begin{align}
H_3 &= \tfrac{1}{2}[\varepsilon_{312}F_{12} + \varepsilon_{321}F_{21}] = \tfrac{1}{2}[(1)F_{12} + (-1)(-F_{12})] = F_{12}.
\end{align}

Summary:
\[
\boxed{H_1 = F_{23}, \qquad H_2 = F_{31}, \qquad H_3 = F_{12}}.
\]

\subsection*{Reverse check: $F_{\alpha\beta} = \varepsilon_{\alpha\beta\mu}H_\mu$}

Compute:
\begin{align}
F_{12} &= \varepsilon_{12\mu}H_\mu = \varepsilon_{123}H_3 = H_3 = F_{12}. \quad \checkmark \\
F_{23} &= \varepsilon_{23\mu}H_\mu = \varepsilon_{231}H_1 = H_1 = F_{23}. \quad \checkmark \\
F_{31} &= \varepsilon_{31\mu}H_\mu = \varepsilon_{312}H_2 = H_2 = F_{31}. \quad \checkmark
\end{align}

And for the antisymmetric partners:
\begin{align}
F_{21} &= \varepsilon_{21\mu}H_\mu = \varepsilon_{213}H_3 = -H_3 = -F_{12}. \quad \checkmark
\end{align}

The correspondence is verified in both directions. \hfill $\blacksquare$

\end{document}
